\documentclass[runningheads]{llncs}
\usepackage{graphicx}
\usepackage{amsmath,amssymb}
\usepackage{booktabs}
\usepackage{color}
\usepackage[width=122mm,left=12mm,paperwidth=146mm,height=193mm,top=12mm,paperheight=217mm]{geometry}
\newcommand{\figlabel}[1]{%
  \special{pdf:dest (#1) [ @thispage /XYZ @xpos @ypos null ]}%
  \label{#1}%
}
\newcommand{\figref}[1]{%
  \special{pdf:bann << /Subtype /Link /Border [0 0 0] /A << /S /GoTo /D (#1) >> >>}%
  Figure~\ref{#1}%
  \special{pdf:eann}%
}
\setcounter{topnumber}{3}
\setcounter{bottomnumber}{2}
\setcounter{totalnumber}{5}
\renewcommand{\topfraction}{0.95}
\renewcommand{\bottomfraction}{0.9}
\renewcommand{\textfraction}{0.05}
\renewcommand{\floatpagefraction}{0.85}
\setlength{\textfloatsep}{8pt plus 2pt minus 2pt}
\raggedbottom

\begin{document}
\pagestyle{headings}
\mainmatter
\def\ECCV16SubNumber{14798}

\title{Comparing Semantic and Logical Composition in Latent Diffusion Models}
\titlerunning{ECCV-16 submission ID \ECCV16SubNumber}
\authorrunning{ECCV-16 submission ID \ECCV16SubNumber}
\author{Molefe Molefe Devon Jarvis Richard Klein}
\institute{14798}

\maketitle

\begin{abstract}
By leveraging the compositional structure of natural language, text-to-image diffusion is able to create coherent scenes with multiple objects while respecting their respective spatial structure. Yet, training these models remains highly data inefficient, as examples of the object pairings that have the desired semantic relationship are needed, motivating the exploration of compositional diffusion techniques. These methods use \emph{logical} rather than \emph{textual} composition by combining specialist diffusion models at inference time to generate complex scenes and aim to generalize to unseen combinations. However, the extent to which the logical composition respects the semantic relationship between objects remains unclear, with examples of ``a cat'' AND ``a dog'' creating chimeras rather than two animals side-by-side. We demonstrate the differences between the two paradigms through a trajectory-level analysis that shows how they diverge early in the diffusion process and favour substantially different endpoints in latent space. Further, we identify the text prompt stable diffusion would need to reproduce a given logical composition, using this to characterize the gap between what logical operators produce and what language can express. Finally, we show that a simple linear interpolation of velocities can provide a semantic grounding of logical composition.
\end{abstract}

% !TEX root = ../../eccv2016submission.tex

\section{Introduction}

Text-to-image diffusion models can produce coherent multi-object scenes from natural-language prompts, but they optimize for caption-level semantic plausibility rather than explicit logical constraints. This matters for compositional generation: prompting Stable Diffusion with ``$A$ and $B$'' and composing single-concept experts with an AND-like operator are not the same objective, even when both are intended to represent conjunction. In practice, the two paradigms can yield qualitatively different outputs, including semantically plausible co-occurrence scenes on one side and hybridized or dominant-mode outcomes on the other.

This paper studies that mismatch directly as a \emph{composability gap}. We compare semantic composition in SD~3.5~\cite{scaling_rectified_flow_transformers_for_high_resolution_image_synthesis} against logical score-space composition from composable diffusion~\cite{liu2022compositional} and SuperDiff AND~\cite{skreta2025the}. To isolate operator effects, all conditions start from the same noise sample per pair and seed. We then analyze both full denoising trajectories and terminal latents, so the gap is characterized not only by final visual differences but also by where and when trajectories separate.

We further test whether AND endpoints are reachable via standard text conditioning. For each SuperDiff-AND output, we project the image back into SD text-conditioning space using SD-IPC~\cite{ding2023sdipc}, a closed-form, training-free projection based on the pseudo-inverse of CLIP's text-projection matrix~\cite{ding2023sdipc}. This is a deliberately conservative choice: SD-IPC is the weakest plausible inverter, so failure here provides a lower-bound argument for structural inaccessibility. If even this minimal projection cannot close the cycle, then AND's target lies outside the region reachable by natural language conditioning.

Our experiments show a consistent pattern: semantic and logical composition diverge under shared-noise initialization, divergence is amplified in late denoising, and the SD-IPC cycle breaks for AND outputs while closing for monolithic text generation. These findings establish the composability gap as a structural property, not a prompting artefact, and open a research programme: understanding its representational cause (why do separately trained experts produce off-manifold compositions?) and designing geometric corrections that enforce manifold consistency during inference.

% !TEX root = ../../eccv2016submission.tex

\section{Related Works}

\noindent\textbf{Text-to-Image Generation.}
Text-to-image (T2I) generation is a conditional generative modeling problem: given a text prompt \(y\), the model samples an image \(x \sim p_\theta(x \mid y)\) that is both visually plausible and semantically aligned with the prompt. T2I lies at the interface of natural language processing and computer vision, and has become a core task in modern AIGC pipelines~\cite{shi2020captioning,ramesh2022hierarchical}.

Recent T2I systems can be grouped into a few practical families. First, hierarchical and autoregressive pipelines model text-conditioned generation via staged decoders (e.g., CLIP-latent priors and image decoders)~\cite{ramesh2022hierarchical}. Second, retrieval-augmented approaches inject external visual neighbors to improve faithfulness and factual grounding~\cite{blattmann2022retrieval,chen2022reimagen}. Third, diffusion-based and flow-based formulations now dominate large-scale generation quality, including latent diffusion variants and transformer backbones~\cite{podell2024sdxl,chen2024pixartsigma}.

In diffusion-style T2I models, the text prompt is encoded into embeddings that condition denoising updates across timesteps. Sampling then iteratively reconstructs an image from noise while preserving prompt semantics. A standard mechanism is classifier-free guidance (CFG), which combines conditional and unconditional predictions to trade off prompt adherence and sample realism~\cite{ho2022classifierfree}. In practice, efficiency-oriented distillation methods can further reduce the required number of denoising steps while retaining competitive quality~\cite{sauer2024add}.

More generally, diffusion models learn generation by reversing a gradual corruption process. A forward Markov process adds noise to real images over many steps, and a learned denoiser approximates the reverse dynamics to map noise back to the data manifold. This iterative formulation is one reason diffusion models are robust for high-dimensional image synthesis: training is stable, conditioning is flexible, and sampling quality scales well with model/data size. Recent systems further improve this paradigm through latent-space modeling and transformer-based rectified-flow formulations~\cite{podell2024sdxl,scaling_rectified_flow_transformers_for_high_resolution_image_synthesis}.

Current state-of-the-art T2I performance is largely achieved by scaled diffusion/flow-transformer systems trained on massive text-image corpora~\cite{podell2024sdxl,chen2024pixartsigma}. Stable Diffusion 3 belongs to this frontier: it uses a rectified-flow transformer formulation and strong text conditioning to improve prompt following, typography, and high-resolution synthesis~\cite{scaling_rectified_flow_transformers_for_high_resolution_image_synthesis}. In this paper, we use SD~3.5 as the semantic composition baseline and compare it against logical score-space composition operators.

\noindent\textbf{Compositional Text-to-Image Generation.}
Compositional text-to-image generation requires a model to satisfy multiple constraints simultaneously in a single image, including object identity, attributes, counts, and inter-object relations. Although modern T2I models produce high-fidelity samples, they still frequently fail on compositional prompts, e.g., by swapping attributes across objects or violating spatial/relational constraints. Prior work has explored compositionality through focused subproblems such as concept conjunction, structured guidance, and attention-based control~\cite{liu2022compositional,kitada-2022-tfsdg,Wu_2023_ICCV}. A complementary evaluation direction is open-world benchmarking: T2I-CompBench shows that compositional robustness remains a key bottleneck even for strong contemporary models~\cite{NEURIPS2023_f8ad010c}. In this paper, we focus on inference-time composition by combining pretrained diffusion experts in score space~\cite{liu2022compositional}.



% !TEX root = ../../eccv2016submission.tex

\section{Method}

We study two mechanisms for combining concepts in Stable Diffusion~\cite{scaling_rectified_flow_transformers_for_high_resolution_image_synthesis}: \emph{semantic composition}, which uses a single natural-language prompt (e.g.\ ``a cat and a dog''), and \emph{logical composition}, which combines concept-specific scores during inference to guide sampling toward a joint high-density region. Our analysis proceeds in three stages. First, we examine shared-noise trajectories to determine when the two mechanisms begin to separate. Next, we analyse hybrid trajectories formed by interpolating their velocity fields. Finally, we search for $p^*$, the natural-language prompt that yields the closest Stable Diffusion reconstruction to SuperDiff AND~\cite{skreta2025the}. Starting from the same seed, we assess whether logical conjunction has an equivalent representation in language space or leaves an irreducible residual gap. Measuring this residual across concept pairs makes the gap directly testable, showing where language-based prompting can match logical conjunction and where the two composition mechanisms consistently diverge. 

All runs use SD~3.5 Medium with the flow-matching ODE sampler, \(50\) denoising steps, and guidance scale \(4.5\). In the main setting, we evaluate \(24\) concept pairs over \(24\) seeds (\(N=576\) pair-seed records per pooled condition) under matched-noise comparisons. For each pair-seed, we run solo prompts, monolithic composition, PoE, SuperDiff-AND, and (when prompt inversion is enabled) prompt-based \(p^*\) reconstructions. We compute metrics per pair-seed record and then pool across all pairs and seeds to obtain condition-level distributions. The reported metrics are trajectory geometry (shared-seed manifold projections), terminal latent gap \(d_T\), temporal accumulation \(\Delta d_t\), and distributional divergences (KL/Jeffrey's), which together capture both magnitude and structure of the composability gap.

\subsection{Semantic, Logical, and Guided Trajectory Dynamics}
\label{sec:trajectory_dynamics}

We compare five conditions of compositional generation, all initiated from a shared Gaussian noise sample $x_T \sim \mathcal{N}(0,I)$. The shared origin is essential: it eliminates stochastic variation as a confound and ensures that any observed difference in trajectories or terminal latents reflects the operator, not the random seed. The five conditions are: (i)~\emph{solo $c_1$} — Stable Diffusion conditioned on the first concept alone; (ii)~\emph{solo $c_2$} — Stable Diffusion conditioned on the second concept alone; (iii)~\emph{semantic composition} — Stable Diffusion conditioned on the single joint natural-language prompt ``$c_1$ and $c_2$''; and \emph{logical composition} in Liu et al.'s composable diffusion framework~\cite{liu2022compositional}, instantiated as (iv)~\emph{PoE} and (v)~\emph{SuperDiff AND}, where SuperDiff combines per-concept velocity fields at every denoising step.

To characterise the geometric relationship between all conditions simultaneously, we record the full latent trajectory ${z_t^{(s)}}_{t=T}$ at every denoising step rather than only the terminal output. We then flatten each latent $z_t \in \mathbb{R}^{C\times H\times W}$ into a vector and fit a single metric MDS to the full matrix of pairwise $\ell_2$ distances across all conditions and timesteps. The joint places every trajectory in a shared coordinate frame, so cross-condition (i.e., between different generation conditions such as solo, semantic, PoE, and SuperDiff) separations are geometrically meaningful rather than artefacts of the projection basis. 

MDS is preferred over PCA because it is distance-centric: PCA finds the 2D subspace of maximum total variance, which may be dominated by within-condition variation along a shared axis; MDS instead directly minimises a stress criterion on pairwise distances, so embedded positions faithfully reflect inter-condition separations in the original $\mathbb{R}^D$ space. Each condition traces a coloured path from the shared origin $x_T$, with colour encoding denoising progress from dark (high noise) to light (terminal). The solo conditions serve as concept attractors: they mark where single-concept flows terminate and provide reference anchors for reading where the compositional conditions land relative to each concept.

\subsection{Hybrid Guidance for Compositional Sampling}
\label{sec:hybrid}

% The semantic and logical paradigms represent extremes of a spectrum: monolithic conditioning is anchored to language but limited by corpus statistics, while SuperDiff AND is logically well-defined but semantically ungrounded. To bridge them, we define a family of \emph{guided hybrid} trajectories by linearly interpolating the velocity fields at each denoising step:
% \begin{equation}
%     v_\alpha = (1-\alpha)\,v_{\mathrm{mono}} + \alpha\,v_{\mathrm{AND}}, \quad \alpha\in[0,1].
%     \label{eq:hybrid}
% \end{equation}
% At $\alpha=0$ this recovers monolithic CFG; at $\alpha=1$ it recovers pure SuperDiff AND. This is not a probabilistic mixture but a geometric interpolation in the velocity space of the flow-matching ODE. Intermediate values trace trajectories that are semantically anchored by the monolithic velocity while being steered toward the joint-density region targeted by AND.

% We additionally consider a time-varying schedule $\alpha(t)$ that decays linearly from $\alpha_{\max}$ to $\alpha_{\min}$. Since global layout and object identity are established in the high-noise regime, setting $\alpha_{\max}=1$ in the first steps drives strong subject separation via AND; decaying toward $\alpha_{\min}$ in later steps allows the monolithic velocity to re-anchor the trajectory semantically, correcting species drift and restoring scene coherence. We sweep $\alpha \in \{0.0, 0.1, 0.3, 0.5, 0.7, 1.0\}$ to characterise the full trade-off.

\subsection{Prompt Design and Closed-Loop Data Generation}
\label{sec:inversion}

We introduce the term \emph{composability gap} to describe the apparent mismatch between semantic and logical composition: starting from the same noise, the monolithic semantic path and the logical AND path can terminate in different latent basins. As a first operational definition, we measure this gap at the terminal step as the distance between the monolithic endpoint and the SuperDiff-AND endpoint.
\begin{equation}
    \Delta_T^{\mathrm{mono}}
    =
    \mathrm{MSE}\!\left(z_T^{\mathrm{mono}}, z_T^{\mathrm{AND}}\right)
    =
    d_T^{\mathrm{mono}}.
    \label{eq:gap}
\end{equation}

This scalar is intentionally a starting probe, not a final theory: it quantifies how far semantic composition is from the logical target before any prompt recovery. This sets up the role of $p^*$ directly: can we find a prompt whose conditioning steers SD3.5 into the same latent basin as SuperDiff AND, thereby reducing (or closing) the apparent gap? If not, the residual indicates a representational limit of language-conditioned generation rather than simple prompt mis-specification.

For image-to-conditioning recovery, we use SD-IPC~\cite{ding2023sdipc}, a closed-form projection that maps a CLIP image embedding back into SD text-conditioning space via the pseudo-inverse of CLIP's text-projection matrix, followed by L2 normalisation and scaling. We treat the resulting conditioning sequence as a surrogate \(p^*\) and feed it directly into SD's cross-attention pathway to synthesize \(x_{\mathrm{SD}}(p^*)\). SD-IPC is training-free and requires no optimisation, making it the most conservative choice: if even this minimal inverter cannot close the regeneration cycle, any more expressive inverter (optimisation-based or learned) can only improve, so the observed cycle gap is a lower bound on structural inaccessibility.

The answer depends on the nature of the composed concepts. When concepts are compatible — as in Liu et al.'s composable diffusion~\cite{liu2022compositional}, where AND generates novel but coherent scenes such as ``a red bear on a tree'' — AND's joint mode may remain on the manifold of natural scene descriptions. In our setting, however, the composed concepts compete at the object level: cat and dog are distinct semantic categories whose joint mode, rather than producing two coexisting animals, tends toward chimeric blending of their feature statistics. 

\textit{A prompt $p^*$ that approximates this target cannot be a better scene description; it must exit the manifold of coherent compositions — something closer to cat-dog chimera'' than to ``a cat and a dog.'' The composability gap is precisely this distance, and we show that the best-approximating natural-language prompt is semantically unnatural, confirming that AND's target lies off the natural scene manifold.}

We demonstrate this through two complementary diagnostics. First, we examine the \emph{semantic character} of each $p^*$ approximation: if AND's target is off the natural scene manifold, caption-based descriptions of its output should become semantically unstable or degenerate rather than clean scene descriptions. Second, we measure the \emph{residual gap} $d_T^{p^*} = \mathrm{MSE}(z_T^{p^*}, z_T^{\mathrm{AND}})$ that remains after reconstruction with training-free natural-language approximations. If this gap persists across complementary prompt-recovery methods, the target is representationally inaccessible to language-conditioned SD3.5 generation.

\input{sections/experiments/experiments}
% !TEX root = ../../eccv2016submission.tex

\section{Conclusion}

This paper introduced and measured a \emph{composability gap} between two notions of conjunction in text-to-image generation: semantic conjunction from monolithic text prompting and logical conjunction from score-space composition (PoE/SuperDiff AND). Under shared-noise initialization, these mechanisms consistently follow different latent trajectories and terminate in different regions, showing that the mismatch is not a decoding artifact but a trajectory-level property of the denoising dynamics.

Across concept pairs and seeds, divergence emerges early but is amplified most strongly in late denoising, where fine-grained identity and attribute conflicts are resolved. This provides a temporal account of the gap: semantic and logical objectives may start from the same stochastic state, yet accumulate most of their separation near the end of sampling.

We further used VLM-based prompt recovery as a counterfactual test of language-space recoverability. BLIP-2 captions used as \(p^*\) move SD outputs closer to SuperDiff AND, but a persistent residual remains in endpoint, trajectory, and pooled distributional analyses. This indicates that the gap is not only prompt-selection error; part of it reflects a structural mismatch between language-conditioned generation and logical score-space intersection.

A practical next step is to design guidance schemes that explicitly trade off semantic naturalness and logical fidelity over time (for example, scheduled or learned hybrid guidance). A broader extension is to test whether the same dynamics hold across other model families, composition operators, and richer multi-concept settings.


\bibliographystyle{splncs}
\bibliography{egbib}
\end{document}
